\section{Dehn filling, variational response, and quantum surgery}
\label{sec:dehnresponse}

There is a second route from volume optimization to a quantum-volume
comparison: perturb a complete cusped structure by prescribing peripheral
holonomy. This route has a particularly explicit answer. The critical-value
response of the real concave optimization recovers the Euclidean cusp
quadratic form. After passage to a holomorphic boundary potential, the
local volume response is exactly the imaginary part of a Legendre
critical value, and gives the filled volume when the geometric completion exists.
Its Taylor coefficients reproduce the Neumann--Zagier expansion, and its
Hessian is the derivative entering the torsion surgery formula.
These are deductions from the established variational, deformation, and
torsion identities cited below; no novelty claim is made.

\subsection{Peripheral conventions and the two small parameters}

Let $M$ be an oriented, finite-volume, one-cusped hyperbolic
three-manifold. Fix an oriented meridian--longitude basis $(\mu,\lambda)$
and the geometric deformation branch through the complete structure.
We use logarithms of the $\mathrm{PSL}_2(\CC)$ multipliers, equivalently
of the eigenvalue ratios, so that
\begin{equation}
 u=H(\mu),\qquad v=H(\lambda),\qquad
 v(u)=\tau u+c_3u^3+c_5u^5+O(u^7),\qquad
 B=\Im\tau>0.
 \label{eq:dehnholonomy}
\end{equation}
Thus $u$ is twice a local logarithm of an $\mathrm{SL}_2$ eigenvalue,
after removing its constant sign at the complete representation.
In particular, the smooth filling equation is $pu+qv=2\pi i$,
with the branch chosen continuously from the cusp. Using eigenvalue
logarithms $\widehat u=u/2$, $\widehat v=v/2$ instead gives
$p\widehat u+q\widehat v=\pi i$ and
$\widehat v=\tau\widehat u+4c_3\widehat u^3+16c_5\widehat u^5+\cdots$.
The oddness and the holomorphic potential are those of
Neumann--Zagier~\cite[Sections 4--6]{NeumannZagier}.

For a primitive slope $(p,q)$ put
\begin{equation}
 A=p+q\tau,\qquad Q(p,q)=\frac{|A|^2}{B},\qquad
 pu+qv(u)=i\theta.
 \label{eq:dehncone}
\end{equation}
Here $\theta$ is an angular holonomy parameter. It is the cone angle
when the local structure has the corresponding geometric cone completion;
$\theta=0$ is the complete cusp and $\theta=2\pi$ gives a smooth
hyperbolic filling when that completion exists. The quantity $Q$ is the squared length of
the slope on a Euclidean cusp cross section divided by its area.
For sufficiently long slopes the small solution stays in one fixed
holonomy neighborhood throughout $0\leq\theta\leq2\pi$. Indeed, put
$h(u)=v(u)-\tau u=O(u^3)$ and $t=q/A$. Since
$|t|\leq B^{-1}$, a sufficiently small disk has
$|t h'(u)|<1/2$ uniformly in the slope. The equation
$u+t h(u)=i\theta/A$ then has a unique small analytic solution,
uniformly for bounded $\theta$ and sufficiently large $|A|$, with
$u=O(|A|^{-1})$. This is an assertion about a holonomy germ;
the hyperbolic Dehn surgery theorem supplies the smooth geometric
completion for sufficiently long primitive slopes
\cite[Section 4]{NeumannZagier}. Geometric completion is a separate
statement from the analytic existence of this small solution.

Two expansions must be kept distinct: $|A|^{-1}\to0$ is a
\emph{classical long-slope expansion}, whereas the quantum level
$r\to\infty$ is a semiclassical expansion. The coefficients $c_3,c_5$
belong to the first. They are not by themselves quantum two-loop
or higher-loop coefficients.

\subsection{What the linear term does to real volume optimization}

Suppose a chosen ideal triangulation of the compact core of $M$ has
a positive complete geometric angle structure, with every ideal vertex
link a torus. Use independent linear coordinates $x$ on its full angle space,
with tetrahedral and edge-angle constraints already imposed, and center
them at the complete maximum $x=0$. Write
\[
 \mathscr V(x)=\sum_{\text{tetrahedra}}\Vol(\Delta_x),\qquad
 V_0=\mathscr V(0),\qquad K=-D^2\mathscr V(0)>0.
\]
Peripheral angular holonomies are affine in the angles and vanish
at $0$ in the selected branches. Thus there is a real linear map $P$
such that
\[
 Px=\binom{\Im u}{\Im v},\qquad
 L=(p\ \ q)P,\qquad Lx=\theta.
\]
The boundary intersection pairing gives $\operatorname{rank}P=2$,
so $L\ne0$. Together with $K>0$, this makes the constrained
stationarity system invertible at $0$ and gives a unique nearby
interior maximizer for sufficiently small $\theta$. These assertions
use the full constrained angle space, not an arbitrarily restricted
family of tetrahedral shapes.

At that maximizer, the leading--trailing edge variations impose the
missing real edge equations; the peripheral variation imposes the
missing real filling equation. This follows from the angle-holonomy
pairing and volume derivative in~\cite[Lemma 4.4, Proposition 4.6 and
Proposition 5.4]{FuterGueritaud}. It first identifies a local edge-glued
holonomy solution. At $\theta=2\pi$, an interior critical point on
the corresponding slice gives a complete hyperbolic filling by their
Theorem 6.1. For other angular slices we retain the local holonomy
interpretation unless a cone completion is specified.
For long slopes the small holonomy solution stays in this positive
geometric neighborhood throughout a bounded angular interval, so the
same constrained-stationarity identification applies there.

To see the peripheral part directly, on the edge-glued deformation
space the volume differential is
\begin{equation}
 \dd V=\frac12\bigl(\Re v\,\dd\Im u
                         -\Re u\,\dd\Im v\bigr).
 \label{eq:dehnvolumedifferential}
\end{equation}
Requiring this covector to be proportional to
$p\,\dd\Im u+q\,\dd\Im v$ gives
$p\Re u+q\Re v=0$. Hence the single angular condition, together
with constrained criticality, gives the full complex equation
\eqref{eq:dehncone}.

The proposed linear perturbation can therefore be written as
\begin{equation}
 \mathscr B(j)=\max_x\{\mathscr V(x)+jLx\},\qquad
 \theta=\mathscr B'(j),\qquad
 V_{p,q}(\theta)=\mathscr B(j)-j\theta.
 \label{eq:realdehnlegendre}
\end{equation}
All statements here are local near the unique interior complete
maximum. The envelope theorem gives
\[
 V_{p,q}'(\theta)=-j.
\]
On a geometric cone branch, the Schl\"afli formula gives
$V_{p,q}'(\theta)=-\ell_{\rm core}/2$,
so the coefficient of the added real linear term is
\begin{equation}
 \boxed{j=\frac{\ell_{\rm core}}2.}
 \label{eq:dehnsourcecore}
\end{equation}
The cone angle is its conjugate variable. In particular, the
critical value $\mathscr B$ of the biased objective and the physical
volume $V_{p,q}$ must be distinguished by the Legendre subtraction
in~\eqref{eq:realdehnlegendre}.

\subsection{The cusp shape in the inverse volume Hessian}

Set
\[
 \chi=L K^{-1}L^T.
\]
Differentiating the stationary equation for the perturbed objective gives
\[
 \left.\frac{\dd x_j}{\dd j}\right|_{j=0}=K^{-1}L^T,
 \qquad \mathscr B''(0)=\chi.
\]
Consequently
\begin{equation}
 V_{p,q}(\theta)=V_0-\frac{\theta^2}{2\chi}+O(\theta^4).
 \label{eq:dehnquadraticresponse}
\end{equation}
Evenness follows from the signed-cone continuation $u\mapsto-u$.
Comparison with the geometric calculation below yields
\begin{equation}
 \boxed{L K^{-1}L^T=\frac{2|p+q\tau|^2}{\Im\tau}=2Q(p,q).}
 \label{eq:cuspHessian}
\end{equation}
Equivalently, writing $a=\Re\tau$,
\begin{equation}
 \boxed{
 P K^{-1}P^T=\frac2B
  \begin{pmatrix}1&a\\a&|\tau|^2\end{pmatrix},\qquad
 \det(P K^{-1}P^T)=4.}
 \label{eq:cuspmatrixHessian}
\end{equation}
Thus the peripheral response of the angle-volume Hessian determines
the cusp shape. Although $K$ and $P$ depend on the triangulation,
this two-by-two matrix depends only on the marked cusp.
The formula assumes the nondegenerate positive angle maximum above;
it is not a claim that an arbitrary chosen ideal triangulation has one.

\subsection{An exact holomorphic critical-value identity}

Define
\begin{equation}
 F(u)=\int_0^u v(t)\,\dd t
     =\frac{\tau u^2}{2}+\frac{c_3u^4}{4}
        +\frac{c_5u^6}{6}+O(u^8).
 \label{eq:dehnprimitive}
\end{equation}
With Neumann--Zagier's convention $\Phi_{\rm NZ}'=2v$,
this is $F=\Phi_{\rm NZ}/2$, after fixing the additive constant.
Assume $q\ne0$ and set
\begin{equation}
 \boxed{
 \mathcal S_{p,q}(u;\theta)
       =F(u)+\frac{p}{2q}u^2-\frac{i\theta}{q}u.}
 \label{eq:dehnsurgeryaction}
\end{equation}
The quadratic term implements the change of peripheral polarization;
the final term is linear in the boundary coordinate. Its critical
equation is exactly $pu+qv=i\theta$. Let $u_\theta$ denote the small
critical solution, and write
$\mathcal S^*(\theta)=\mathcal S_{p,q}(u_\theta;\theta)$.

\begin{proposition}[Local filling critical value]
\label{prop:dehncriticalvalue}
On the local geometric deformation branch through the complete cusp,
with $V_{p,q}$ interpreted as its volume germ (or the volume of its
geometric completion when this exists),
\begin{equation}
 \boxed{V_{p,q}(\theta)-V_0=\frac12\Im\mathcal S^*(\theta).}
 \label{eq:dehnexactcriticalvolume}
\end{equation}
Moreover,
\begin{equation}
 \frac{\dd\mathcal S^*}{\dd\theta}=-\frac{i u_\theta}{q},\qquad
 \frac{\dd^2\mathcal S^*}{\dd\theta^2}
       =\frac{1}{q\bigl(p+qv'(u_\theta)\bigr)}.
 \label{eq:dehnenvelope}
\end{equation}
\end{proposition}

\begin{proof}
Choose integers $\rho,\sigma$ with $p\sigma-q\rho=1$.
The complementary peripheral holonomy is
\[
 w=\rho u+\sigma v=\frac{i\sigma\theta-u}{q}.
\]
When a geometric cone completion exists, in these orientation conventions
the physical core complex length
has positive real part when represented by $-w$, and therefore
$\ell_{\rm core}=\Re(u/q)$ near the positive cone branch.
The envelope theorem gives the first identity
in~\eqref{eq:dehnenvelope}. Its half imaginary part is
$-\ell_{\rm core}/2$, exactly the Schl\"afli derivative of volume.
Both sides of~\eqref{eq:dehnexactcriticalvolume} vanish at $\theta=0$.
Differentiating $pu_\theta+qv(u_\theta)=i\theta$ gives the second
identity in~\eqref{eq:dehnenvelope}.

There is also an algebraic proof valid directly on the local holonomy
branch. Neumann--Zagier's volume formula is
\begin{align}
 V(u)-V_0
 &=\frac14\Im(u\bar v)
     +\Im\left(\frac12F(u)-\frac14uv\right)\notag\\
 &=\frac12\Im F(u)-\frac12\Re u\,\Im v.
 \label{eq:NZrealvolume}
\end{align}
Substituting $\Im v=(\theta-p\Im u)/q$ yields
\eqref{eq:dehnexactcriticalvolume}. This proof avoids assuming
an independently constructed cone path.
\end{proof}

This is an exact identity of analytic germs, so it compares all
orders of the two classical expansions. For $q=0$, exchange the
oriented basis to $(\mu',\lambda')=(\lambda,-\mu)$; the new slope
has $q'=-p\ne0$. The real angle optimization itself needs no such
basis change. Primitive $q=0$ slopes have $p=\pm1$ and do not
occur in a long-slope sequence in a fixed basis.
Figure~\ref{fig:mono-dehn-filling} records the local intersection
condition and the response at the complete structure.

% BEGIN INLINE FIGURE: 04_dehn_filling.tex
\begin{figure}[htbp]
\centering
\begin{tikzpicture}[font=\small,line cap=round,>=Latex]
  \begin{scope}[shift={(-5.6,-1.1)}]
    \draw[->,mpgray!70] (-.2,0)--(3.7,0) node[below] {$u$};
    \draw[->,mpgray!70] (0,-.3)--(0,3.0) node[left] {$v$};
    \draw[mpblue,line width=1.2pt,domain=-.1:3.1,samples=80]
      plot (\x,{.55*\x+.045*\x*\x*\x});
    \draw[mpochre,line width=1pt] (.05,2.22682)--(3.3,-.21068);
    \fill[mpblue] (0,0) circle (2pt) node[below left=3pt,align=right] {complete\\structure};
    \fill[mpochre] (1.6,1.06432) circle (2.5pt);
    \draw[mpgray!65,dashed] (1.6,0)--(1.6,1.06432)--(0,1.06432);
    \node[above left=2pt,text=mpblue] at (2.9,2.697) {$v=v(u)$};
    \node[above right=2pt,text=mpochre] at (1.6,1.06432) {$u_\theta$};
    \node[text=mpgray,font=\scriptsize] at (1.8,3.4) {schematic projection of complex curves};
  \end{scope}
  \node[align=left,text width=7.0cm,anchor=west] at (-.65,.25) {%
    \textcolor{mpochre}{$\displaystyle pu+qv=i\theta$}\\[7pt]
    $\displaystyle F'(u)=v(u)$\\[7pt]
    $\displaystyle\mathcal S_{p,q}(u;\theta)
       =F(u)+\frac{p}{2q}u^2-\frac{i\theta}{q}u\quad(q\ne0)$\\[7pt]
    $\displaystyle\partial_u\mathcal S_{p,q}=0
       \quad\Longleftrightarrow\quad pu+qv(u)=i\theta$\\[10pt]
    \textbf{Response at the complete structure}\\[5pt]
    $\displaystyle\chi=LK^{-1}L^T
       =\frac{2|p+q\tau|^2}{\Im\tau}$.};
\end{tikzpicture}
\caption[Dehn filling in peripheral phase space]{Local Dehn filling as an intersection in the complex peripheral
holonomy space. The drawing records the intersection condition; it is
not a real slice of the complex deformation curve. Here $u,v$ are
logarithms of multipliers, $\theta=0$ is the complete structure, and
$\theta=2\pi$ gives smooth filling when the branch has that geometric
completion. The Hessian response formula assumes the
nondegenerate positive angle maximum specified in the text.}
\label{fig:mono-dehn-filling}
\end{figure}
% END INLINE FIGURE

\subsection{The Neumann--Zagier coefficients through sixth order}

Series inversion of~\eqref{eq:dehncone} gives
\begin{equation}
 u_\theta=\frac{i\theta}{A}
   +\frac{iqc_3\theta^3}{A^4}
   +i\theta^5\left(\frac{3q^2c_3^2}{A^7}
                         -\frac{qc_5}{A^6}\right)+\cdots.
 \label{eq:dehninversion}
\end{equation}
Substitution in the action, or integration of the envelope derivative,
gives
\begin{equation}
 \mathcal S^*(\theta)=\frac{\theta^2}{2qA}
  +\frac{c_3\theta^4}{4A^4}
  +\frac{\theta^6}{6}
       \left(\frac{3qc_3^2}{A^7}-\frac{c_5}{A^6}\right)+\cdots.
 \label{eq:dehncriticalseries}
\end{equation}
Taking half its imaginary part yields the volume loss
\begin{align}
 V_0-V_{p,q}(\theta)
 ={}&\frac{B\theta^2}{4|A|^2}
       -\frac{\theta^4}{8}\Im\frac{c_3}{A^4}\notag\\
    &-\frac{\theta^6}{12}
       \Im\left(\frac{3qc_3^2}{A^7}-\frac{c_5}{A^6}\right)
       +O(|A|^{-8}).
 \label{eq:dehnconeloss}
\end{align}
Here the displayed remainder is a long-slope estimate for bounded
$\theta$, on the common holonomy neighborhood constructed above.
The uniform analytic inverse expands in $i\theta/A$ with coefficients
analytic in the bounded parameter $q/A$; the next volume terms are
therefore $O(|A|^{-8})$. Since $q/A$ is bounded, the three displayed terms have
orders $|A|^{-2},|A|^{-4},|A|^{-6}$. For a fixed slope and
$\theta\to0$, the corresponding Taylor remainder is $O(\theta^8)$.

For the smooth hyperbolic filling $M_{p/q}$ this becomes
\begin{align}
 \boxed{\begin{aligned}
 V_0-\Vol(M_{p/q})
 ={}&\frac{\pi^2B}{|A|^2}
       -2\pi^4\Im\frac{c_3}{A^4}\\
    &-\frac{16\pi^6}{3}
       \Im\left(\frac{3qc_3^2}{A^7}-\frac{c_5}{A^6}\right)
       +O(|A|^{-8}).
 \end{aligned}}
 \label{eq:dehnsmoothloss}
\end{align}
These are the Neumann--Zagier coefficients, with the multiplier
normalization fixed explicitly~\cite[Section 6]{NeumannZagier}.
In particular the leading loss is $\pi^2/Q(p,q)$, and the quartic
term detects the first nonlinear deformation coefficient $c_3$.

\subsection{Computing the quartic coefficient from angle-volume derivatives}

The comparison supplies a more detailed computational identity.
Keep $K,L,\chi$ as above, and write
\[
 T_3=D^3\mathscr V(0),\qquad T_4=D^4\mathscr V(0),\qquad
 e=\frac{K^{-1}L^T}{\chi},\qquad
 R=K^{-1}-\frac{K^{-1}L^TLK^{-1}}{\chi}.
\]
Let $t=T_3(e,e,\cdot)$ be the resulting column covector. Solving
the constrained stationary equation to second order gives
\[
 x_\theta=\theta e+\frac{\theta^2}{2}Rt+O(\theta^3).
\]
Indeed $Le=1$, $LR=0$, and $RK$ is the identity on
$\ker L$. Inserting this expression into the Taylor expansion of
$\mathscr V$ gives
\begin{align}
 V_{p,q}(\theta)={}&V_0-\frac{\theta^2}{2\chi}
  +\frac{\theta^3}{6}T_3[e,e,e]\notag\\
  &+\frac{\theta^4}{24}
       \left(T_4[e,e,e,e]+3t^TRt\right)+O(\theta^5).
 \label{eq:dehnquartictensors}
\end{align}
The cross term with the unknown cubic displacement vanishes because
that displacement is in $\ker L$ and $Ke=L^T/\chi$.
The remaining quadratic-displacement terms combine to $t^TRt/8$.
Evenness of the geometric branch then removes the odd coefficients
and improves the remainder to $O(\theta^6)$. Comparison with
\eqref{eq:dehnconeloss} yields
\begin{equation}
 \boxed{T_3[e,e,e]=0,\qquad
 T_4[e,e,e,e]+3t^TRt
      =3\Im\frac{c_3}{(p+q\tau)^4}.}
 \label{eq:dehnquarticdictionary}
\end{equation}
Thus $c_3$ can be accessed from derivatives of the angle-volume
functional by varying the slope. The correction $3t^TRt$ accounts
for readjustment of all other angles. Keeping only the raw fourth
derivative would miss that response.

\subsection{The closed quantum-volume statement and its local consequence}

A precise closed-manifold analogue is the Chen--Yang volume conjecture.
At the specified root $\zeta_r=e^{2\pi i/r}$, with $r$ odd, its
absolute-growth statements are
\begin{equation}
 \frac{4\pi}{r}\log|\mathrm{RT}_r(N;\zeta_r)|\longrightarrow\Vol(N),
 \qquad
 \frac{2\pi}{r}\log\mathrm{TV}_r(N;\zeta_r)\longrightarrow\Vol(N).
 \label{eq:closedquantumvolume}
\end{equation}
The root and normalization matter; these are the Chen--Yang versions
for hyperbolic manifolds~\cite{ChenYang}. We use absolute growth
to avoid mixing the different Chern--Simons sign conventions in
the literature.

For a boundary state whose geometric semiclassical phase satisfies
\[
 \Phi(u)=\Phi_0+\frac12F(u),\qquad \Im\Phi_0=V_0,
\]
the filling phase is
\begin{equation}
 \mathcal T_\theta(u)=\Phi_0+\frac12\mathcal S_{p,q}(u;\theta).
 \label{eq:dehnquantumphase}
\end{equation}
Possible real framing constants do not alter its imaginary part.
With $\hbar_r=4\pi i/r$, Proposition~\ref{prop:dehncriticalvalue}
immediately gives
\begin{equation}
 \left|e^{\mathcal T_\theta(u_\theta)/\hbar_r}\right|
   =\exp\left(\frac{r}{4\pi}V_{p,q}(\theta)\right).
 \label{eq:dehnquantumexponent}
\end{equation}
In particular, the deformation of the geometric exponential is
\begin{equation}
 \log\left|\frac{e^{\mathcal T_\theta(u_\theta)/\hbar_r}}
                         {e^{\Phi_0/\hbar_r}}\right|
       =-\frac{r}{4\pi}\bigl(V_0-V_{p,q}(\theta)\bigr).
 \label{eq:dehnquantumresponse}
\end{equation}
Every classical coefficient in~\eqref{eq:dehnconeloss} therefore
appears with precisely the required quantum exponential coefficient.
At $\theta=2\pi$, Yoshida's complex-volume formula has this
normalization; see the unnumbered filling formula and the potential
normalization (5.9) on~\cite[p.~456]{WongYangDehn}.
Quadratic and linear filling kernels also appear explicitly
in~\cite[equations (12)--(14)]{GangRomoYamazaki}; their boundary
coordinates and level conventions must be converted before numerical
coefficients are compared.

\begin{proposition}[Conditional local quantum transfer]
\label{prop:dehnquantumconditional}
Suppose an exact surgery representation for a specified quantum
invariant has, near the geometric branch, phase
$\mathcal T_\theta$ and a uniform classical-symbol expansion of
its amplitude. Suppose the geometric critical point is nondegenerate,
its steepest-descent cycle occurs with nonzero coefficient in the
specified integration cycle, the net leading stationary-phase
coefficient of that geometric sector is uniformly bounded away
from zero, and the remaining integral or sum is
exponentially smaller, uniformly on a chosen compact deformation
interval. Then stationary phase gives the volume exponent
\eqref{eq:dehnquantumexponent} for the full invariant, with a
uniform expansion on that interval. Without the last global
condition, the same statement holds for the specified local
geometric contribution.
\end{proposition}

\begin{proof}
At the geometric critical point,
\[
 \mathcal T_\theta''(u_\theta)
   =\frac{p+qv'(u_\theta)}{2q}\ne0.
\]
Uniform stationary phase on the stipulated cycle gives a nonzero
amplitude times $e^{\mathcal T_\theta(u_\theta)/\hbar_r}$,
with powers of $r^{-1}$ and a Gaussian power of $r$ according
to the normalization of the kernel. Such powers contribute only
$O(\log r/r)$ to normalized logarithmic growth. The exponential
gap permits the addition of the remaining terms without changing
the leading growth. The critical-value proposition identifies the
exponent with the geometric volume.
\end{proof}

For a general knot this proposition identifies the needed analytic
input; it does not establish that input from a complete-point
Kashaev asymptotic alone. A value at $u=0$ does not provide a
uniform boundary-state expansion near $0$, and a local critical
point does not by itself specify its coefficient in the global
surgery sum. The regular tetrahedral Airy formula in
Section~\ref{sec:quantumAiry} concerns a different integral and
does not supply that missing manifold-level information.

\subsection{The Gaussian factor and the torsion surgery formula}

The derivative of the filling map is
\begin{equation}
 D_*:=p+qv'(u_\theta),\qquad
 \frac{\dd u_\theta}{\dd\theta}=\frac{i}{D_*},\qquad
 \mathcal S_{uu}(u_\theta;\theta)=\frac{D_*}{q}.
 \label{eq:dehnjacobian}
\end{equation}
Thus the same quantity controls the classical susceptibility and
the Gaussian surgery integral. With the usual square-root
normalization of the boundary Fourier kernel, its contribution
has dependence $D_*^{-1/2}$, up to the common normalization and
phase choices. A full amplitude also contains the boundary-state
torsion and a solid-torus factor.

More precisely, suppose $\theta=2\pi$ gives a smooth hyperbolic filled
manifold and the representation is its holonomy. Assume its adjoint
representation is acyclic and its restriction to the compact core is
regular for both $\mu$ and $p\mu+q\lambda$: the peripheral invariant
vector paired with either curve supplies a basis of the one-dimensional
twisted $H_1$, with the compatible boundary basis in $H_2$.
In the local holonomy chart this requires $D_*\ne0$; the core is
loxodromic, so $\sinh(w/2)\ne0$. The change-of-peripheral-curve
and surgery formulas for adjoint Reidemeister torsion then give
\begin{align}
 \operatorname{Tor}_{p\mu+q\lambda}(M_u)
   &=\bigl(p+qv'(u)\bigr)\operatorname{Tor}_\mu(M_u),\notag\\
 \operatorname{Tor}(M_{p/q})
   &=\frac{\operatorname{Tor}_{p\mu+q\lambda}(M_u)}
               {4\sinh^2(w/2)},\qquad
 w=\rho u+\sigma v.
 \label{eq:dehntorsionsurgery}
\end{align}
These are understood up to the standard torsion signs and
consistent choices of square root. They are the identities
used in~\cite[equations (5.16) and (5.23)]{WongYangDehn}.
For the general regularity and change-of-curve statements, see
\cite[Propositions 5.13--5.14]{PortiTorsionSurvey}; its
Proposition 4.23(b), specialized to the adjoint representation
$(k_1,k_2)=(2,0)$, gives the surgery factor up to the torsion sign.
Consequently the expected full geometric amplitude is proportional to
\begin{equation}
 \boxed{
 \frac{2\sinh(w/2)}
      {\sqrt{\bigl(p+qv'(u)\bigr)\operatorname{Tor}_\mu(M_u)}}
       =\frac{1}{\sqrt{\operatorname{Tor}(M_{p/q})}}.}
 \label{eq:dehngaussiantorsion}
\end{equation}
The identity shows what must accompany the boundary Hessian.
Matching the Gaussian factor alone would leave out both the cusp
torsion and the solid-torus amplitude. This one-loop comparison
uses more information than the classical volume exponent.
The closed-manifold torsion in this formula is evaluated at smooth
filling; it is not being assigned to a generic cone holonomy.
At the cusp endpoint $\theta=0$ the solid-torus factor vanishes,
so a uniform quantum statement including that endpoint requires
a separate normalization and analysis. The compact parameter
interval in Proposition~\ref{prop:dehnquantumconditional} is
subject to its nonvanishing-amplitude hypothesis.

\subsection{A figure-eight benchmark and the scope of the conclusion}

For the figure-eight knot complement, in the standard peripheral
basis Neumann--Zagier compute~\cite[equation (68)]{NeumannZagier}
\begin{equation}
 v(u)=2\sqrt3\,i\left(u+\frac{u^3}{3}
                    +\frac{23u^5}{180}+O(u^7)\right).
 \label{eq:figureeightNZcoefficients}
\end{equation}
Thus $\tau=2\sqrt3\,i$, $c_3=2i/\sqrt3$, and
$c_5=23i/(30\sqrt3)$. The first two volume-loss terms are
\begin{align}
 V_0-\Vol(M_{p/q})={}&\frac{2\sqrt3\,\pi^2}{p^2+12q^2}\notag\\
 &-\frac{4\pi^4}{\sqrt3}
  \frac{p^4-72p^2q^2+144q^4}{(p^2+12q^2)^4}
  +O\bigl((p^2+q^2)^{-3}\bigr).
 \label{eq:figureeightfillingloss}
\end{align}
This agrees directly with the expansion following equation (68)
in~\cite{NeumannZagier}. The angle-Hessian response is correspondingly
\[
 \chi=\frac{p^2+12q^2}{\sqrt3}.
\]

An independent calculation starts with the two regular ideal tetrahedra
and their Lobachevsky volume. In the Neumann--Zagier presentation their
shapes $z_1,z_2$ satisfy
\[
 z_1z_2(z_1-1)(z_2-1)=1,\qquad
 u=\log\bigl(z_2(1-z_1)\bigr),\qquad
 v=2\log\bigl(z_1(1-z_1)\bigr).
\]
Eliminating the shapes from these equations gives the local peripheral
relation
\[
 \cosh(v/2)=1+\cosh u-\cosh(2u).
\]
Equivalently, with the eigenvalue variables
$m=e^{u/2}$ and $l=-e^{v/2}$, it is
$l+l^{-1}=m^4-m^2-2-m^{-2}+m^{-4}$, the nonabelian
figure-eight $A$-polynomial relation in this choice of signs.
The sign in $l$ is essential: its complete value is $-1$ whereas
$v(0)=0$. The equation and the chosen tangent $v'(0)=2\sqrt3\,i$
recover the cubic and quintic coefficients above. This is a local
branch calculation, not a classification of the global curve or of fillings.

Use $\alpha_i=\arg z_i$, $\beta_i=-\arg(1-z_i)$, and
$\gamma_i=\pi-\alpha_i-\beta_i$. Choose the independent variations
$\xi=(\delta\alpha_1,\delta\beta_1,\delta\alpha_2)$.
The edge equation then makes the six angle variations
\[
 E\xi=(\xi_1,\xi_2,-\xi_1-\xi_2,
       \xi_3,\xi_1-\xi_2+\xi_3,-\xi_1+\xi_2-2\xi_3).
\]
At the complete structure all six angles are $\pi/3$.
Since $\Lambda''(\alpha)=-\cot\alpha$, direct differentiation gives
\begin{equation}
 K=\frac1{\sqrt3}
 \begin{pmatrix}4&-1&3\\-1&4&-3\\3&-3&6\end{pmatrix},\qquad
 P=\begin{pmatrix}0&-1&1\\2&-2&0\end{pmatrix},\qquad
 P K^{-1}P^T=\frac1{\sqrt3}
 \begin{pmatrix}1&0\\0&12\end{pmatrix}.
 \label{eq:figureeightdirectHessian}
\end{equation}
This recovers the cusp shape directly from the optimizer.
Put $D_0=p^2+12q^2$. The unit angular-response vector is
\[
 e=\frac{(6q-p,-6q-p,p-6q)^T}{2D_0}.
\]
Using
\[
 \Lambda'''(\pi/3)=\frac43,\qquad
 \Lambda''''(\pi/3)=-\frac8{3\sqrt3},
\]
the six entries of $Ee$ occur in opposite pairs, so
$T_3[e,e,e]=0$. Direct contraction of the fourth derivative and
the constrained inverse-Hessian correction gives
\begin{equation}
 T_4[e,e,e,e]+3t^TRt
   =\frac{2\sqrt3\,(p^4-72p^2q^2+144q^4)}{D_0^4}.
 \label{eq:figureeightdirectquartic}
\end{equation}
Multiplying by $\theta^4/24$ and setting $\theta=2\pi$ gives
the quartic filled-volume term in~\eqref{eq:figureeightfillingloss}
with the opposite sign, as required for volume rather than volume loss.
Thus the real optimization and the holomorphic series agree in a
calculation made directly from the tetrahedral volume derivatives.

For this example there is also a theorem about the actual quantum
invariants: Wong--Yang prove the leading RT asymptotic with relative
error $O(r^{-1})$, and the
RT/TV volume conjectures for every hyperbolic rational Dehn filling
of the figure-eight knot~\cite[Theorems 1.1 and 1.2]{WongYangDehn}.
Their leading amplitude is the inverse square root of adjoint torsion,
with their specified overall constant. Thus the local consistency
calculation above is contained in a stronger result in that example.
The present discussion reorganizes the comparison around volume
optimization and its derivatives; it does not produce a new case
of the quantum-volume conjecture.

For a general knot, the unconditional results here are the classical
critical-value identity and its coefficient dictionary. The quantum
conclusion for a specified local geometric contribution follows under
the stated stationary-phase hypotheses. Passing to the full invariant
requires control of its exact surgery representation and all competing
contributions. One should first keep the filling slope fixed while
taking $r\to\infty$, and then study the long-slope coefficients.
A simultaneous limit in which the slope grows with $r$ requires
uniform estimates not supplied by separate fixed-slope expansions.

The computational companion in
\path{experiments/quantum_surgery/README.md}, with output at
\path{artifacts/quantum_surgery/result.json}, checks finite formal
series, canonical-coordinate identities, and the figure-eight
angle-Hessian contractions in exact arithmetic. These fixtures
check the displayed coefficient algebra. They do not compute a
filled hyperbolic manifold or an exact quantum invariant, and do
not establish convergence, a global integration cycle, or its
noncancellation.

\subsection{The symplectic pairing behind the calculation}
\label{sec:NZsymplectic}

The pairing in Neumann--Zagier is the structure behind the preceding
formulas. It identifies the peripheral variables as conjugate coordinates,
makes the geometric deformation locus Lagrangian, and fixes the canonical
transformation for Dehn filling. We record the normalization explicitly,
since interchanging meridian and longitude or using eigenvalue rather
than multiplier logarithms changes the displayed signs or factors.

\subsubsection{From the tetrahedra to the boundary intersection form}

For $n$ ideal tetrahedra introduce independent ambient coordinates
\[
 Z_j=\log z_j,\qquad W_j=\log(1-z_j),\qquad
 \Omega_{\rm tet}=\sum_{j=1}^n\dd Z_j\wedge\dd W_j.
\]
The actual shapes satisfy $e^{Z_j}+e^{W_j}=1$. Their product is a
complex Lagrangian of dimension $n$ in this $2n$-dimensional phase space.
The logarithmic edge equations and peripheral holonomies are affine
linear functions of $X=(Z_1,\ldots,Z_n,W_1,\ldots,W_n)^T$.

Let $C_\lambda,C_\mu$ be the longitude and meridian coefficient rows,
let $R_E$ be the edge-equation matrix, and set
\[
 J_{2n}=\begin{pmatrix}0&I_n\\-I_n&0\end{pmatrix},\qquad
 U=\begin{pmatrix}C_\lambda\\C_\mu\\R_E\end{pmatrix}.
\]
For a triangulated manifold with $h$ torus cusps, the original matrix identity is
\begin{equation}
 \boxed{U J_{2n}U^T
 =2\begin{pmatrix}0&I_h&0\\-I_h&0&0\\0&0&0\end{pmatrix}.}
 \label{eq:NZmatrixpairing}
\end{equation}
This is~\cite[Theorem 2.2, equation (21)]{NeumannZagier}, in their
longitude--meridian row order. They normalize the coefficient-row
pairing as $\langle a,b\rangle=\tfrac12aJ_{2n}b^T$.
Equivalently, the unnormalized coefficient pairing of peripheral curves
is twice their oriented intersection number; edge rows pair to zero
with one another and with all peripheral rows. The topological
interpretation is proved in their Section 3. A geometric construction
of symplectic bases is given in~\cite{MathewsPurcellSymplectic}.

The vanishing edge-edge pairings say that the linear constraints
Poisson-commute. Since $\operatorname{rank}R_E=n-h$, their
coisotropic reduction has complex dimension $2h$. Corollary 2.4 of
\cite{NeumannZagier} identifies the reduced form as
\begin{equation}
 \boxed{\Omega_\partial=\frac12\sum_{i=1}^h
                   \dd v_i\wedge\dd u_i,\qquad
 \alpha_\partial=\frac12\sum_{i=1}^h v_i\,\dd u_i,
 \qquad \dd\alpha_\partial=\Omega_\partial.}
 \label{eq:NZboundarysymplectic}
\end{equation}
Locally near the regular geometric point, imposing the edge equations
on the tetrahedral Lagrangian and reducing gives the geometric
deformation Lagrangian $\mathcal L_M$ in boundary holonomy space.
For a modern formulation of the reduction, see~\cite{DimofteVeen}.
The symplectic form lives on the ambient boundary phase space;
its pullback to $\mathcal L_M$ is zero.

\subsubsection{The potential and the volume primitive}

Writing $v=v(u)$ on the geometric branch, Lagrangianity gives
\begin{equation}
 \sum_i\dd v_i\wedge\dd u_i=0
 \quad\Longleftrightarrow\quad
 \frac{\partial v_i}{\partial u_j}
       =\frac{\partial v_j}{\partial u_i}.
 \label{eq:NZreciprocity}
\end{equation}
Thus $\sum_i v_i\,\dd u_i=\dd F$ locally, and
$\alpha_\partial|_{\mathcal L_M}=\dd(F/2)$.
This is the reciprocity and potential statement of
\cite[Theorem 3]{NeumannZagier}. For one cusp the existence of a
primitive is automatic once $v=v(u)$; the substantial content is
the triangulation-level pairing and, for several cusps, the
symmetry of the mixed derivatives. In the original normalization,
$\Phi_{\rm NZ}=2F$, whereas the complex-volume phase used above is
$\Phi=\Phi_0+F/2$.

For one cusp write
\[
 u=x+iy,\qquad v=r+is.
\]
The real volume one-form is an exact modification of the imaginary
part of the symplectic primitive:
\begin{equation}
 \beta=\frac12(r\,\dd y-x\,\dd s)
       =\Im\alpha_\partial-\frac12\dd(xs),\qquad
 \dd\beta=\Im\Omega_\partial.
 \label{eq:NZvolumeprimitive}
\end{equation}
It is therefore closed on $\mathcal L_M$. Integrating it gives exactly
\eqref{eq:NZrealvolume}. The factor $1/2$ can also be seen from
the leading--trailing angle variations: angular holonomy varies by
twice the boundary intersection number, while volume varies by the
real holonomy~\cite[Lemma 4.4 and Proposition 5.4]{FuterGueritaud}.
The same pairing thus supplies the real optimization differential
and the holomorphic generating function.

\subsubsection{Filling as a canonical transformation}

Choose $\rho,\sigma\in\ZZ$ with $p\sigma-q\rho=1$ and use
the filling slope as the new meridian:
\begin{equation}
 a=pu+qv,\qquad b=\rho u+\sigma v,\qquad
 \frac12\dd b\wedge\dd a=\frac12\dd v\wedge\dd u.
 \label{eq:NZcanonicalfilling}
\end{equation}
For $q\ne0$, the type-I generating function is
\begin{equation}
 \boxed{G(u,a)=\frac{p}{4q}u^2-\frac{ua}{2q}
                         +\frac{\sigma}{4q}a^2.}
 \label{eq:NZcanonicalkernel}
\end{equation}
Indeed, $v=(a-pu)/q$ and $b=(\sigma a-u)/q$ give the exact identity
\begin{equation}
 \dd G=\frac12b\,\dd a-\frac12v\,\dd u.
 \label{eq:NZgeneratingidentity}
\end{equation}
Thus
\[
 \widetilde\Phi(a)=\operatorname{crit}_u\bigl(\Phi(u)+G(u,a)\bigr)
 \quad\text{satisfies}\quad
 \widetilde\Phi'(a)=\frac b2.
\]
The stationary equation is $a=pu+qv$. Setting $a=i\theta$ gives
\begin{equation}
 \widetilde\Phi(i\theta)
 =\Phi_0+\frac12\mathcal S^*_{p,q}(\theta)
          -\frac{\sigma\theta^2}{4q}.
 \label{eq:NZcanonicalaction}
\end{equation}
The last term is real, hence invisible to the volume. At
$\theta=2\pi$ it is $-\pi^2\sigma/q$, the framing term in
the Yoshida normalization used in~\cite[p.~456, formula preceding (5.9)]{WongYangDehn}.
The quadratic and linear terms of the surgery action are therefore
fixed by a symplectic change of coordinates.

The complementary core holonomy $b$ is conjugate to the filling
holonomy $a$. Along $a=i\theta$,
\[
 \dd\Im\widetilde\Phi
   =\Im\left(\frac12b\,i\dd\theta\right)
   =\frac12\Re b\,\dd\theta
   =-\frac12\ell_{\rm core}\,\dd\theta.
\]
This gives Schl\"afli and the real source
$j=\ell_{\rm core}/2$ in the canonical coordinates.

On the canonical relation, $G=(ab-uv)/4$, so
\begin{equation}
 \boxed{\widetilde\Phi-\frac{ab}{4}
                  =\Phi-\frac{uv}{4}.}
 \label{eq:NZinvariantpotential}
\end{equation}
Equivalently $\widetilde\Phi_{\rm NZ}-ab=\Phi_{\rm NZ}-uv$,
where $\widetilde\Phi_{\rm NZ}=4(\widetilde\Phi-\Phi_0)$.
The equality is local with consistently fixed logarithmic branches
and additive constants. It explains why the combination of potential
and peripheral product occurs in the complex-volume formula.
The differential of the boxed expression is the invariant one-form
$\tfrac14(v\,\dd u-u\,\dd v)$.

\subsubsection{Covariance of the filling coefficients}

The powers of $p+q\tau$ in~\eqref{eq:dehnsmoothloss} also have
a canonical-coordinate interpretation. Write the same geometric
Lagrangian in the new coordinates as
\[
 b(a)=\widetilde\tau a+\widetilde c_3a^3
                         +\widetilde c_5a^5+O(a^7).
\]
Using the inverse series of $a=pu+qv(u)$ gives, with $A=p+q\tau$,
\begin{align}
 \widetilde\tau&=\frac{\rho+\sigma\tau}{p+q\tau},
 &\Im\widetilde\tau&=\frac{\Im\tau}{|p+q\tau|^2},\notag\\
 \widetilde c_3&=\frac{c_3}{A^4},
 &\widetilde c_5&=\frac{c_5}{A^6}-\frac{3qc_3^2}{A^7}.
 \label{eq:NZcoefficientcovariance}
\end{align}
The cancellation in the cubic coefficient is
\[
 \sigma A-q(\rho+\sigma\tau)=1.
\]
It is exactly the symplectic determinant condition. Let
$\widetilde F(0)=0$ and $\widetilde F'=b$.
Since $\Re a=0$ on the cone locus, the real boundary correction
in the new coordinates vanishes. Therefore
\begin{equation}
 V_{p,q}(\theta)-V_0
  =\frac12\Im\widetilde F(i\theta)
  =-\frac{\theta^2}{4}\Im\widetilde\tau
    +\frac{\theta^4}{8}\Im\widetilde c_3
    -\frac{\theta^6}{12}\Im\widetilde c_5+\cdots.
 \label{eq:NZvolumenewpolarization}
\end{equation}
This is the previously computed volume expansion in filling
coordinates. The cubic coefficient has a weight-$-4$ covariance
law under a peripheral basis change; this alone does not assert
that it is a modular form as a function of cusp shape.
The pairing fixes the transformation laws, while the actual values
of $c_3,c_5,\ldots$ describe nonlinear geometry of $\mathcal L_M$.

\subsubsection{An exact Hessian and Monge--Ampere identity}

There is a useful strengthening of~\eqref{eq:cuspmatrixHessian}.
Eliminate internal angles at fixed peripheral angular holonomies,
writing
\[
 \overline V(y,s)=\max_{P\xi=(y,s)}\mathscr V(\xi).
\]
All statements are local in the positive geometric neighborhood,
where the maximum is interior and nondegenerate. Let
$K(\xi)=-D^2\mathscr V(\xi)>0$ at the maximizing point.
The boundary intersection pairing gives $\operatorname{rank}P=2$,
so the constrained implicit-function theorem supplies a unique nearby
interior maximizer for every sufficiently small $(y,s)$.
Stationarity in the internal directions enforces real edge gluing,
so this maximizer lies on the same deformation branch $v=v(u)$.
The map $u=x+iy\mapsto(y,\Im v(u))$ is locally invertible because
$\Im v'(0)=B>0$. It therefore identifies the angular boundary
parameters with that geometric branch, rather than with an unrelated
constrained family.
The Schur-complement identity for the linear constraints is
\begin{equation}
 -D^2\overline V=(P K(\xi)^{-1}P^T)^{-1}.
 \label{eq:NZreducedHessian}
\end{equation}
Set $T=v'(u)=A_0+iB_0$, with $B_0>0$ near the complete point.
The volume differential and Cauchy--Riemann equations give
\[
 \overline V_y=\frac r2,\qquad
 \overline V_s=-\frac x2,\qquad
 \dd s=B_0\dd x+A_0\dd y,\qquad
 \dd r=A_0\dd x-B_0\dd y.
\]
Solving these two linear differential equations yields
\begin{equation}
 \boxed{-D^2\overline V
 =\frac1{2B_0}\begin{pmatrix}|T|^2&-A_0\\-A_0&1\end{pmatrix},
 \qquad \det(-D^2\overline V)=\frac14.}
 \label{eq:NZMongeAmpere}
\end{equation}
Equivalently,
\begin{equation}
 P K(\xi)^{-1}P^T
  =\frac2{\Im v'(u)}
       \begin{pmatrix}1&\Re v'(u)\\
                       \Re v'(u)&|v'(u)|^2\end{pmatrix}.
 \label{eq:NZnonlinearresponse}
\end{equation}
The determinant-$4$ identity holds throughout this local branch.
At $u=0$ it reduces to the cusp formula already derived.
For a filling path it gives the exact response
\begin{equation}
 V_{p,q}''(\theta)
   =-\frac{\Im v'(u_\theta)}
            {2|p+qv'(u_\theta)|^2},
 \label{eq:NZexactcurvature}
\end{equation}
in agreement with~\eqref{eq:dehnenvelope}.

For a geometric interpretation of the positive Hessian metric
$g=-D^2\overline V$, multiplication by $i$ in the $u$-coordinate
induces, in coordinates $(y,s)$,
\[
 J_T=\frac1{B_0}
      \begin{pmatrix}-A_0&1\\-|T|^2&A_0\end{pmatrix},
 \qquad J_T^2=-I.
\]
Then
\begin{equation}
 \omega_{\rm ang}=-\frac12\dd y\wedge\dd s,\qquad
 g(X,Y)=\omega_{\rm ang}(X,J_TY),\qquad
 g=\frac{B_0}{2}(\dd x^2+\dd y^2).
 \label{eq:NZcompatiblemetric}
\end{equation}
This real area form comes from the angular polarization and the
positive complex tangent. It is not the pullback of
$\Omega_\partial$, which vanishes on the geometric Lagrangian.
Nor has it been identified with the earlier
Thurston--Deligne--Mostow complex hyperbolic metric.

For completeness, assume the same positive complete triangulation and
regular deformation hypotheses with $h$ torus cusps, and use all $2h$
peripheral angular constraints. Their matrix $P$ has rank $2h$ by the
boundary intersection pairing. Let
$T=(\partial v_i/\partial u_j)=A_0+iB_0$.
It is symmetric, and $B_0>0$ near the complete point. In coordinates
$(y,s)=(\Im u,\Im v)$ the same calculation gives
\begin{equation}
 -D^2\overline V
  =\frac12\begin{pmatrix}
   B_0+A_0B_0^{-1}A_0&-A_0B_0^{-1}\\
   -B_0^{-1}A_0&B_0^{-1}
   \end{pmatrix},\qquad
 \det(-D^2\overline V)=4^{-h}.
 \label{eq:NZmulticuspHessian}
\end{equation}
At the complete structure $T$ is diagonal, with the individual cusp
shapes on the diagonal. Coupling between distinct cusps appears in
higher derivatives of the even Neumann--Zagier potential, for example
through terms $u_i^2u_j^2$.

\subsubsection{The quantum Gaussian as a half-density transformation}

Let a geometric boundary state have leading form
$A_0(u)e^{\Phi(u)/\hbar}$, with $\Phi'=v/2$.
A quantization of the canonical transformation uses the exponential
kernel $e^{G(u,a)/\hbar}$ and its mixed-derivative normalization.
The relevant derivatives are
\[
 -G_{ua}=\frac1{2q},\qquad
 (\Phi+G)_{uu}=\frac{p+qv'(u)}{2q}.
\]
The kernel normalization divided by the stationary Gaussian therefore
has the dependence
\begin{equation}
 \boxed{\sqrt{\frac{-G_{ua}}{(\Phi+G)_{uu}}}
    =\frac1{\sqrt{p+qv'(u)}}
    =\sqrt{\frac{\dd u}{\dd a}}.}
 \label{eq:NZhalfdensity}
\end{equation}
Universal powers of $\hbar$, phases, and square-root branches are
fixed by the chosen integral convention. The displayed dependence is
the transformation law for a half-density along $\mathcal L_M$.
It is exactly the law obeyed by the inverse square root of peripheral
torsion in~\eqref{eq:dehntorsionsurgery}. The extra
$2\sinh(b/2)$ factor comes from the solid torus and is additional
input. Explicit quantum filling kernels appear
in~\cite[equations (12)--(14)]{GangRomoYamazaki}.

The symplectic structure fixes the conjugate variables, the canonical
generating function, and this Gaussian Jacobian. It does not select
the integration cycle or prove that the geometric contribution
survives cancellation in the exact quantum invariant. The analytic
hypotheses in Proposition~\ref{prop:dehnquantumconditional} remain
necessary for that conclusion.


